\pdfoutput=1
\documentclass[preprints,article,submit,oneauthor]{Definitions/mdpi}
\usepackage{mathtools}
\usepackage{mathrsfs}
\setlength{\headheight}{18pt}
\numberwithin{equation}{section}

\newcommand{\C}{\mathbb C}
\newcommand{\R}{\mathbb R}
\newcommand{\D}{\mathbb D}
\newcommand{\spec}{\operatorname{spec}}
\newcommand{\spr}{\operatorname{spr}}
\newcommand{\adj}{\operatorname{adj}}
\newcommand{\conv}{\operatorname{conv}}
\newcommand{\Rea}{\operatorname{Re}}
\newcommand{\Ima}{\operatorname{Im}}
\newcommand{\Int}{\operatorname{int}}

\firstpage{1}
\makeatletter
\setcounter{page}{\@firstpage}
\makeatother
\pubvolume{1}
\issuenum{1}
\articlenumber{0}
\pubyear{2026}
\copyrightyear{2026}
\datereceived{}
\daterevised{}
\dateaccepted{}
\datepublished{}

\Title{Rigidity at the Extremal Crouzeix Constant}
\newcommand{\orcidauthorA}{0009-0009-5346-753X}
\Author{Shanmu Jin\orcidA{}}
\AuthorNames{Shanmu Jin}
\address{Department of Neurosurgery, Peking Union Medical College
Hospital, No.~1 Shuaifuyuan, Dongcheng District, Beijing 100730, China\\
Email: \href{mailto:jinshanmu1129@pku.edu.cn}{jinshanmu1129@pku.edu.cn}}
\corres{Correspondence: \href{mailto:jinshanmu1129@pku.edu.cn}
{jinshanmu1129@pku.edu.cn}}
\abstract{For a complex matrix \texorpdfstring{\(A\)}{A}, let \texorpdfstring{\(\psi(A)\)}{psi(A)} be the norm of the polynomial functional calculus on its numerical range. We prove that \texorpdfstring{\(\psi(A)=2\)}{psi(A)=2} only when the numerical range is a nondegenerate closed disk. The proof first extracts the equality conditions in the sharp double-layer estimate. After boundary eigenvalues are removed, an extremizer obtained from a finite Blaschke product satisfies an exact boundary-kernel identity. On a curved exposed arc this identity makes the extremizer rational. A direct resolvent argument then places every zero of its reduced numerator inside the numerical range and every pole outside; consequently its full unit lemniscate coincides with the boundary of the numerical range. The defining polynomial of the projective dual of that algebraic boundary divides the Hermitian Kippenhahn determinant. Hyperbolicity and strict convexity force the dual curve to have degree two, and its points at infinity force the resulting conic to be a circle.}
\keyword{numerical range; Crouzeix constant; spectral set; rational lemniscate; hyperbolic polynomial}
\MSC{47A12; 47A25; 15A60}

\begin{document}

\section{Introduction}

Throughout, \(\C\) and \(\R\) denote the complex and real fields,
\(M_N(\C)\) the algebra of
\(N\times N\) complex matrices, and \(I\) the identity matrix of the size
dictated by context.  For a matrix \(B\), the symbols \(B^*\), \(\spec(B)\),
and \(\spr(B)\) denote its conjugate transpose, spectrum, and spectral radius.
The notation \(\|\cdot\|\) denotes the Hilbert-space norm on vectors and the
induced operator norm on bounded linear maps; on \(\C^N\) these are the
Euclidean and induced Euclidean norms.  We put
\[
 \Rea B=\frac{B+B^*}{2},\qquad \Ima B=\frac{B-B^*}{2i}.
\]
For scalars and functions, \(\Rea\) and \(\Ima\) have their usual pointwise
meanings.  For Hermitian matrices
\(B,C\), the relation \(B\succeq C\) means that \(B-C\) is positive
semidefinite.  We write \(\Int E\) for the interior of a planar set \(E\) and
\(\D=\{z\in\C:|z|<1\}\) for the open unit disk.  Scalar matrix coefficients
are written as \(x^*By=\langle By,x\rangle\), where inner products are linear
in the first argument.  For a bounded scalar-valued function \(g\) on a set
\(E\), write \(\|g\|_E=\sup_{z\in E}|g(z)|\).

For \(A\in M_N(\C)\), set
\[
 W(A)=\{u^*Au:u\in\C^N,\ \|u\|=1\}.
\]
For a polynomial \(p\), put
\[
 \psi(A)=
 \sup\left\{\|p(A)\|:p\in\C[z],
               \|p\|_{W(A)}\leq1\right\}.
\]
For a compact convex body \(K\), let \(C(K)\) denote the complex-valued
continuous functions on \(K\), and write
\[
 \mathcal A(K)=\{f\in C(K):f|_{\Int K}\text{ is holomorphic}\}.
\]
The space \(H^\infty(\Int K)\) consists of the bounded holomorphic functions
on \(\Int K\) and carries the norm \(\|\cdot\|_{\Int K}\).  Whenever
\(\spec(A)\subset\Int K\), the notation \(f(A)\) denotes the standard
finite-dimensional holomorphic functional calculus, equivalently evaluation
of any Hermite interpolating polynomial with the required spectral jets.
The constant \(\psi(A)\) was introduced by Crouzeix~\cite{Crouzeix2004}, who
proved \(\psi(A)\leq2\) for \(2\times2\) matrices and conjectured the same
bound in every dimension.  He subsequently obtained the first
dimension-free bound~\cite{Crouzeix2007}.  The disk case is the classical
Okubo--Ando theorem~\cite{OkuboAndo1975}, while the positive operator-valued
double-layer representation behind the general convex-domain estimates goes
back to Delyon and Delyon~\cite{DelyonDelyon1999}.  Crouzeix and Palencia
reduced the universal constant to \(1+\sqrt2\)~\cite{CrouzeixPalencia2017}.

Jin~\cite{Jin2026} and, independently, Lorist and
Schwenninger~\cite{LoristSchwenninger2026} proved the scalar conjecture.  The
latter proof applies an abstract \(2\)-dilation perturbation lemma
simultaneously to all iterates in the double-layer representation.
Subsequently, {\AA}hag, Czy\.z, Per\"al\"a, and Virtanen established sharp
square-function inequalities, the complete conjecture in dimension three,
and equality, stability, and representing-measure
results~\cite{AhagEtAl2026}.  Their equality theory also produces a
two-dimensional reducing nilpotent block and strong dimension-three
restrictions.

This paper develops the equality data into an all-dimensional geometric
rigidity theorem.  Its central steps are rational collapse of the extremizer,
identification of the full lemniscate with the numerical-range boundary, and
a hyperbolic projective-dual contact count.

\begin{Theorem}\label{thm:main}
For every \(N\geq1\) and \(A\in M_N(\C)\),
\[
 \psi(A)=2\quad\Longrightarrow\quad
 W(A)=\{z\in\C:|z-\gamma|\leq\rho\}
\]
for some \(\gamma\in\C\) and some \(\rho>0\).
\end{Theorem}

The proof has four steps.  An equality analysis of the sharp bound produces
two extremal vectors and a pointwise boundary-kernel identity.  A curved
exposed arc then turns that identity into a rational formula for the
extremizer.  A direct resolvent argument identifies its full unit lemniscate
with \(\partial W(A)\).  Finally, the Kippenhahn determinant contains the dual
curve as a hyperbolic factor.  A generic contact count makes the dual a conic,
and the circular points at infinity then force \(\partial W(A)\) to be a
circle.  A separate ingredient, strict domain monotonicity of the matrix
\(H^\infty\) functional-calculus norm, rules out a proper disk produced by the
boundary-spectrum induction.
Extensive numerical optimization by Greenbaum and
Overton~\cite{GreenbaumOverton2018} repeatedly found stationary
configurations with ratio \(2\) whose numerical ranges were disks; the theorem
explains that pattern at exact equality.

\section{Numerical-range preliminaries}

Here and below, \(\conv E\) denotes the convex hull of \(E\).
For a Hermitian matrix \(B\), \(\lambda_{\max}(B)\) denotes its largest
eigenvalue.  At a differentiability point \(\sigma\) of a convex boundary,
\(\mathfrak n(\sigma)\) denotes its outward unit normal; \(\nu\), with or
without a subscript, denotes an arbitrary supporting unit normal.

\begin{Lemma}[Numerical-range geometry]\label{lem:foundation}
The following facts hold.
\begin{enumerate}
\item \(W(A)\) is nonempty, compact, and convex.
\item If \(W(A)\) has empty interior, then \(A\) is normal and
\(\psi(A)=1\).
\item For \(\alpha\in\C\setminus\{0\}\), \(\beta\in\C\), and unitary
\(\mathcal U\),
\[
 W(\alpha A+\beta I)=\alpha W(A)+\beta,\qquad
 \psi(\alpha A+\beta I)=\psi(A),\qquad
 \psi(\mathcal U^*A\mathcal U)=\psi(A).
\]
Also \(W(A^*)=\overline{W(A)}\) and \(\psi(A^*)=\psi(A)\).
\item For a finite nonempty direct sum,
\[
 W\left(\bigoplus_jA_j\right)=\conv\Bigl(\bigcup_jW(A_j)\Bigr),
 \qquad
 \psi\left(\bigoplus_jA_j\right)\leq\max_j\psi(A_j).
\]
\item The support function of \(W(A)\) is
\[
 h_A(\theta)=\lambda_{\max}\!\left(\Rea(e^{-i\theta}A)\right).
\]
\end{enumerate}
\end{Lemma}

\begin{proof}
The unit sphere is compact and \(u\mapsto u^*Au\) is continuous, proving
nonemptiness and compactness.  To prove convexity, take two points of \(W(A)\)
and compress \(A\) to the span of corresponding unit vectors.  If the points
are distinct, those vectors are linearly independent; if they coincide, there
is no segment to prove.  It is therefore enough to know that the numerical
range of a \(2\times2\) matrix is convex.  After unitary
triangularization such a matrix has the form
\(
 \begin{psmallmatrix}\lambda_1&\tau\\0&\lambda_2\end{psmallmatrix}
\).
For \(u=(\cos t,e^{i\varphi}\sin t)\), where
\(0\leq t\leq\pi/2\) and \(\varphi\in\R\),
\[
 u^*Au=\lambda_1\cos^2t+\lambda_2\sin^2t
       +\tau e^{i\varphi}\sin t\cos t.
\]
These points fill the possibly degenerate ellipse with foci
\(\lambda_1,\lambda_2\), hence a convex set.  The two-dimensional compression
is contained in \(W(A)\), so the segment joining the original two points is
contained in \(W(A)\).

If \(W(A)\) has empty interior, its convexity places it in an affine line.
If \(W(A)=\{\gamma\}\), then \(x^*(A-\gamma I)x=0\) for all \(x\).  The complex
polarization identity makes every matrix coefficient of \(A-\gamma I\) zero.
If \(W(A)\) lies in a line, a translation and a nonzero complex rotation make
all quadratic forms of \(A\) real.  The skew-Hermitian part then has zero
quadratic form and vanishes by polarization, so the transformed matrix is
Hermitian.  A normal matrix satisfies
\[
 \|p(A)\|=\|p\|_{\spec(A)}\leq\|p\|_{W(A)}.
\]
The constant polynomial \(p\equiv1\) gives the reverse inequality.

The affine and unitary formulas follow by changing variables in polynomials.
For adjoints use \(\widetilde p(z)=\overline{p(\overline z)}\), for which
\(\widetilde p(A^*)=p(A)^*\).  For a direct sum, a unit vector decomposes into block
components and its numerical value is a convex combination of block numerical
values.  Also
\(p(\bigoplus_jA_j)=\bigoplus_jp(A_j)\); the supremum norm on the convex hull
dominates that on every block.  Finally,
\[
 \max_{z\in W(A)}\Rea(e^{-i\theta}z)
 =\max_{\|u\|=1}u^*\Rea(e^{-i\theta}A)u,
\]
which is the largest eigenvalue of the displayed Hermitian matrix.
\end{proof}

\section{An equality-preserving sharp-bound argument}

We use the equality-preserving form of the Lorist--Schwenninger iterate
argument~\cite{LoristSchwenninger2026}, which records the pointwise identities
needed for rigidity.

\begin{Lemma}[An abstract dilation estimate]\label{lem:dilation}
Let \(T\) act on a nonzero finite-dimensional Hilbert space.  Suppose that \(\mathcal M\) is a
contraction on another Hilbert space, \(\mathcal V\) is an isometry, and
\[
 E_n=2\mathcal V^*\mathcal M^{*n}\mathcal V-T^{*n}\qquad(n\geq1)
\]
are uniformly bounded and commute with \(T\).  Then \(\|T\|\leq2\).

If additionally \(\|T\|=2\), \(\spr(T)<1\), \(T^*Tx=4x\), \(\|x\|=1\), and
\(y=Tx/2\), then
\[
 T^*x=0,\quad T^*y=2x,\quad
 \mathcal M\mathcal Vx=\mathcal Vy,\quad
 \mathcal M^*\mathcal Vy=\mathcal Vx.
\]
\end{Lemma}

\begin{proof}
Put \(\kappa=\|T\|\).  If \(\kappa\leq1\), the conclusion is immediate; assume
\(\kappa>1\).  Choose a unit vector \(x\) with \(T^*Tx=\kappa^2x\), and set
\[
 m_n=\Rea\bigl(x^*E_nT^nx\bigr),\qquad
 \omega=\mathcal M^*\mathcal VTx-\kappa\mathcal Vx.
\]
The defining identity gives
\(\|T^{*n}\|\leq2+\sup_k\|E_k\|\), so the powers of \(T\), and hence
\((m_n)\), are bounded.  Put
\(Q_n=\mathcal V^*\mathcal M^{*n}\mathcal V\) and
\(q_n=T^{*n}x\).  Commutativity of \(E_n\) with \(T\) gives
\[
 2(Q_{n+1}T-TQ_{n+1})
 =T^{*(n+1)}T-TT^{*(n+1)}.
\]
Moreover,
\(m_n=2\Rea(q_n^*Q_nx)-\|q_n\|^2\) and
\(
 \mathcal V^*\mathcal M^{*n}\omega=Q_{n+1}Tx-\kappa Q_nx.
\)
Multiplying the displayed commutator identity on the left by \(q_n^*\), using
\(T^*Tx=\kappa^2x\), and taking real parts now gives
\begin{align*}
 \kappa m_n-m_{n+1}
 &=\kappa(\kappa-1)\|T^{*n}x\|^2
   -2\Rea\bigl((\mathcal V^*\mathcal M^{*n}\omega)^*T^{*n}x\bigr)\\
 &=\kappa(\kappa-1)
   \left\|T^{*n}x-\frac{\mathcal V^*\mathcal M^{*n}\omega}
   {\kappa(\kappa-1)}\right\|^2
   -\frac{\|\mathcal V^*\mathcal M^{*n}\omega\|^2}{\kappa(\kappa-1)}\\
 &\geq-\frac{\|\omega\|^2}{\kappa(\kappa-1)}.
\end{align*}
Telescoping, with the boundedness of \(m_n\), yields
\[
 \kappa m_1
 =\sum_{n\geq1}\kappa^{-n+1}(\kappa m_n-m_{n+1})
 \geq-\frac{\|\omega\|^2}{(\kappa-1)^2}.               \tag{3.1}
\]
On the other hand, using that \(\mathcal V\) is isometric and \(\mathcal M\) contractive,
\begin{align*}
 \|\omega\|^2
 &\leq2\kappa^2
 -2\kappa\Rea\bigl(x^*\mathcal V^*\mathcal M^*\mathcal VTx\bigr)\\
 &=2\kappa^2-\kappa m_1-\kappa^3.                     \tag{3.2}
\end{align*}
Combining (3.1) and (3.2) gives
\[
 \|\omega\|^2\left(1-\frac1{(\kappa-1)^2}\right)
 \leq\kappa^2(2-\kappa),
\]
which is impossible for \(\kappa>2\).

Now let \(\kappa=2\).  Equations (3.1)--(3.2) become
\(2m_1\geq-\|\omega\|^2\) and \(\|\omega\|^2\leq-2m_1\), so equality holds in both.
For each \(n\), the slack in the estimate preceding (3.1) is
\[
 2\left\|T^{*n}x-\tfrac12\mathcal V^*\mathcal M^{*n}\omega\right\|^2
 +\tfrac12\bigl(\|\omega\|^2-\|\mathcal V^*\mathcal M^{*n}\omega\|^2\bigr).
\]
Their positively weighted sum is zero.
Thus, for every \(n\),
\[
 \mathcal V^*\mathcal M^{*n}\omega=2T^{*n}x,
 \qquad
 \|\mathcal V^*\mathcal M^{*n}\omega\|=\|\omega\|.
\]
Since \(\spr(T)<1\), \(T^{*n}x\to0\), and these two equalities force
\(\omega=0\).  Their first instance then gives \(T^*x=0\).
With \(y=Tx/2\), \(\mathcal M^*\mathcal Vy=\mathcal Vx\).  Both
\(\mathcal Vx\) and \(\mathcal Vy\) are unit vectors, so equality in
Cauchy--Schwarz for \(\mathcal M\) gives \(\mathcal M\mathcal Vx=\mathcal Vy\).
Finally \(T^*y=T^*Tx/2=2x\).
\end{proof}

\begin{Proposition}[The sharp numerical-range estimate]\label{prop:sharp}
For every \(N\geq1\), every \(A\in M_N(\C)\), and every polynomial \(p\),
\[
 \|p(A)\|\leq2\|p\|_{W(A)}.                          \tag{3.3}
\]
More generally, if \(K\) is a compact convex body,
\(W(A)\subset K\), \(\spec(A)\subset\Int K\), and \(f\in\mathcal A(K)\), then
\[
 \|f(A)\|\leq2\|f\|_K.                               \tag{3.4}
\]
For \(f\in H^\infty(\Int K)\), the corresponding estimate is
\[
 \|f(A)\|\leq2\|f\|_{\Int K}.
\]
\end{Proposition}

\begin{proof}
Normalize \(\|f\|_K\leq1\).  Cauchy's boundary formula below also holds when
\(f\) is only in \(\mathcal A(K)\): for \(z_0\in\Int K\), the functions
\(f(z_0+r(z-z_0))\), \(r<1\), are holomorphic on a neighborhood of \(K\) and
converge uniformly to \(f\); apply the usual formula to them and pass to the
limit.  The positive operator-valued double-layer kernel used below originates
in Delyon and Delyon~\cite{DelyonDelyon1999} and appears in the
Cauchy-transform form of Crouzeix and Palencia~\cite{CrouzeixPalencia2017};
see also the recent double-layer and configuration-constant treatments
in~\cite{SchwenningerDeVries2025,MalmanEtAl2025}.  Parametrize the
rectifiable convex Jordan curve
\(\Gamma=\partial K\) counterclockwise by arclength.  Its outward unit normal
\(\mathfrak n(\sigma)\) exists almost everywhere and has a measurable
representative.  All boundary identities below are understood
arclength-almost everywhere.  Put
\(R_\sigma=(\sigma I-A)^{-1}\) and
\[
 \mathcal P(\sigma)=\frac1\pi\Rea\bigl(\mathfrak n(\sigma)R_\sigma\bigr).
\]
At a point with a normal, the supporting-half-plane property gives
\[
 \Delta_\sigma:=\Rea\bigl(\overline{\mathfrak n(\sigma)}(\sigma I-A)\bigr)\succeq0.
\]
Multiplication by the two resolvents shows
\[
 \Rea\bigl(\mathfrak n(\sigma)R_\sigma\bigr)
 =R_\sigma^*\Delta_\sigma R_\sigma\succeq0.          \tag{3.5}
\]
Since \(d\sigma=i\mathfrak n(\sigma)\,ds\) and
\(d\overline\sigma=-i\overline{\mathfrak n(\sigma)}\,ds\),
\[
 \mathcal P(\sigma)\,ds
 =\frac1{2\pi i}R_\sigma\,d\sigma
  -\frac1{2\pi i}R_\sigma^*\,d\overline\sigma.
\]
Cauchy's formula gives \(\int_\Gamma R_\sigma\,d\sigma=2\pi iI\); taking
adjoints gives
\(\int_\Gamma R_\sigma^*\,d\overline\sigma=-2\pi iI\).  Therefore
\[
 \int_\Gamma \mathcal P(\sigma)\,ds=2I.              \tag{3.6}
\]
The matrix field \(\mathcal P\) is measurable and positive semidefinite almost
everywhere.  Since the positive-square-root map is continuous on the cone of
positive semidefinite matrices, \(\sigma\mapsto\mathcal P(\sigma)^{1/2}\) is
measurable (with arbitrary values on the exceptional null set).  Consequently
\[
 (\mathcal Vx)(\sigma)=2^{-1/2}\mathcal P(\sigma)^{1/2}x
\]
defines an isometry into \(L^2(\Gamma;\C^N)\).  Let \(M_f\) be multiplication
by \(f\), and set \(T=f(A)\).  For every \(n\geq1\), Cauchy's formula and
the preceding decomposition give
\[
 2\mathcal V^*M_f^{*n}\mathcal V
 =\frac1{2\pi i}\int_\Gamma
   \overline{f(\sigma)^n}R_\sigma\,d\sigma+T^{*n},
\]
because
\[
 \int_\Gamma\overline{f(\sigma)^n}R_\sigma^*\,d\overline\sigma
 =\left(\int_\Gamma f(\sigma)^nR_\sigma\,d\sigma\right)^*
 =-2\pi iT^{*n}.
\]
Thus
\[
 E_n:=2\mathcal V^*M_f^{*n}\mathcal V-T^{*n}
 =\frac1{2\pi i}\int_\Gamma
   \overline{f(\sigma)^n}R_\sigma\,d\sigma.           \tag{3.7}
\]
The right side commutes with \(T\), because every resolvent does.  Also
\[
 \sup_n\|E_n\|\leq
 \frac{\operatorname{length}(\Gamma)}{2\pi}
 \max_{\sigma\in\Gamma}\|R_\sigma\|<\infty.
\]
Lemma~\ref{lem:dilation} proves (3.4).

For \(f\in H^\infty(\Int K)\), normalize
\(\|f\|_{\Int K}\leq1\), choose a Riemann map
\(\phi:\Int K\to\D\), and put \(g=f\circ\phi^{-1}\).  For \(0<r<1\), define
\[
 f_r(z)=g(r\phi(z)).
\]
The Carath\'eodory extension of \(\phi\) shows that \(f_r\in\mathcal A(K)\),
and \(\|f_r\|_K\leq1\).  Moreover, \(g(r\,\cdot)\to g\) locally uniformly with
all derivatives.  The finitely many spectral jets therefore converge, so
\(f_r(A)\to f(A)\); applying (3.4) and passing to the limit proves the
\(H^\infty\) estimate.

For an arbitrary matrix \(A\), apply (3.4) to the outer parallel convex body
\[
 K_\varepsilon=W(A)+\varepsilon\overline\D,
\]
which contains \(W(A)\) in its interior.  For every polynomial \(p\), uniform
continuity on a fixed compact neighborhood gives
\(\|p\|_{K_\varepsilon}\to\|p\|_{W(A)}\).  Letting
\(\varepsilon\downarrow0\) proves (3.3), including the boundary-spectrum and
nonsmooth cases.
\end{proof}

\begin{Proposition}[Equality data]\label{prop:equality-data}
Let \(K=W(A)\) have nonempty interior, let \(\spec(A)\subset\Omega:=\Int K\),
and let \(f\in\mathcal A(K)\) be nonconstant with
\[
 \|f\|_K=1,\qquad |f|=1\ \text{on }\partial K,\qquad \|f(A)\|=2.
\]
Then there are orthonormal vectors \(x,y\) such that, with \(T=f(A)\),
\[
 Tx=2y,\qquad T^*x=0,\qquad T^*y=2x,                 \tag{3.8}
\]
and, for almost every \(\sigma\in\partial K\),
\[
 \mathcal P(\sigma)^{1/2}\bigl(y-f(\sigma)x\bigr)=0.  \tag{3.9}
\]
\end{Proposition}

\begin{proof}
Apply the construction in Proposition~\ref{prop:sharp}.  The spectral mapping
theorem gives \(\spr(T)<1\), because the finitely many eigenvalues of \(A\) are
inside \(\Omega\) and a nonconstant inner function has modulus strictly less
than \(1\) there.  Choose \(x\) with \(T^*Tx=4x\), and put \(y=Tx/2\).
Lemma~\ref{lem:dilation} gives \(T^*x=0\), \(T^*y=2x\), and
\(M_f\mathcal Vx=\mathcal Vy\).  Here \(\|y\|=1\), and
\(2x^*y=x^*Tx=(T^*x)^*x=0\), so \(x,y\) are orthonormal.
The multiplication identity is exactly (3.9).
\end{proof}


\begin{Lemma}[Boundary eigenvalues reduce]\label{lem:boundary-eigen}
If \(Ax=\lambda x\), \(\|x\|=1\), and \(\lambda\in\partial W(A)\), then
\(A^*x=\overline\lambda x\).  Consequently, with \(\simeq\) denoting unitary
similarity,
\[
 A\simeq A_{\partial}\oplus A_{\mathrm{in}},
\]
where \(A_{\partial}\) is diagonal, with the boundary eigenvalues as its
diagonal entries, while
\(\spec(A_{\mathrm{in}})\subset\Int W(A)\).
\end{Lemma}

\begin{proof}
Choose a supporting direction \(\nu\), \(|\nu|=1\), at \(\lambda\).  Equality in
the Rayleigh quotient gives
\[
 \Rea(\overline\nu A)x=\Rea(\overline\nu\lambda)x.
\]
Substituting \(Ax=\lambda x\) gives \(A^*x=\overline\lambda x\).  Thus
\(\C x\) reduces \(A\).  Split it off and repeat.  Every eigenvalue belongs to
\(W(A)\), since an eigenvector realizes it, so the remaining spectrum is in
the interior.
\end{proof}

\begin{Lemma}[Conformal and finite-interpolation facts]\label{lem:complexfacts}
Let \(\Omega\) be a bounded convex domain.
\begin{enumerate}
\item There is a conformal bijection \(\phi:\Omega\to\D\), and it extends to a
homeomorphism \(\overline\Omega\to\overline\D\).
\item At each of finitely many distinct points of \(\Omega\), prescribe a
contiguous jet, consisting of the value and all derivatives through a fixed
finite order.  If these data are realized by a function of
\(H^\infty(\Omega)\) of norm at most \(1\), they are also realized by
\(\mathcal B\circ\phi\), where \(\mathcal B\) is a finite Blaschke product.  If
the infimum of the supremum norms over all interpolants of the prescribed data
is \(1\), then the Schur-class interpolant is unique and is of this form.
\end{enumerate}
\end{Lemma}

\begin{proof}
A bounded convex domain is a Jordan domain.  The Riemann mapping theorem and
the Carath\'eodory boundary theorem therefore give the first assertion.

Transfer the jet data to \(\D\) by \(\phi^{-1}\).  The confluent
change-of-variables matrix between the original and transferred jets is
triangular with nonzero diagonal, because \((\phi^{-1})'\neq0\).  Write the
distinct disk nodes as \(\zeta_i\), prescribe derivatives through order
\(m_i-1\), and let \(q\) be the Hermite polynomial with those jets.  For a norm
bound \(\rho>0\), the generalized Pick matrix used here is indexed by
\((i,r)\), \(0\leq r<m_i\), and has entries
\[
 P(\rho)_{(i,r),(j,s)}=
 \left.
 \frac{1}{r!s!}\,
 \partial_z^r\partial_{\overline w}^{s}
 \frac{\rho^2-q(z)\overline{q(w)}}{1-z\overline w}
 \right|_{z=\zeta_i,\,w=\zeta_j}.                  \tag{3.10}
\]
This matrix depends only on the prescribed jets, not on the chosen Hermite
representative.  The finite confluent Nevanlinna--Pick theorem says that the
data have an interpolant of norm at most \(\rho\) exactly when
\(P(\rho)\succeq0\); see
Sarason~\cite{Sarason1967} and
Bultheel--Lasarow~\cite{BultheelLasarow2008}.  At the norm bound
\(\rho=1\), a positive definite matrix admits a Schur parametrization; choosing
a unimodular terminal parameter gives a rational inner solution.  In the scalar
disk case, rational inner functions are finite Blaschke products.  For general
\(\rho\), the corresponding conclusion is \(\rho\) times a rational inner
function, obtained by rescaling the data.  At the minimal norm, singularity
yields the following uniqueness conclusion.

Contiguous jets at distinct nodes impose independent
conditions: Hermite interpolation supplies polynomials \(h_{i,r}\) whose
normalized derivative functionals satisfy
\(h_{i,r}^{(s)}(\zeta_j)/s!=\delta_{ij}\delta_{rs}\).  Thus there are no hidden
linear dependencies in (3.10).  Now let \(\mathcal J\) send a holomorphic
function on \(\Omega\) to the prescribed jet vector \(\mathbf d\), and put
\[
 c_0=\inf\{\|g\|_\Omega:g\in H^\infty(\Omega),\ \mathcal Jg=\mathbf d\}.
\]
If \(c_0=1\), then \(P(1)\) is singular: if it were positive definite,
continuity of (3.10) in \(\rho\) would give \(P(\rho)\succeq0\) for some
\(\rho<1\), contradicting minimality.  Let
\[
 \Theta(z)=\prod_i
 \left(\frac{z-\zeta_i}{1-\overline{\zeta_i}z}\right)^{m_i}.
\]
Two bounded holomorphic functions have the same prescribed jets exactly when
their difference belongs to \(\Theta H^\infty\).  Thus the interpolation data
form one coset of \(H^\infty/\Theta H^\infty\), represented on the
finite-dimensional model space
\(K_\Theta=H^2\ominus\Theta H^2\) by an operator \(X\) commuting with the
compressed shift.  Sarason's lifting theorem identifies the least interpolant
norm with \(\|X\|\), so \(\|X\|=1\).  The norm of \(X\) is attained because
\(K_\Theta\) is finite dimensional.  Sarason's maximal-vector uniqueness
theorem~\cite[Proposition~5.1]{Sarason1967} now gives a unique norm-one
interpolant and says that it is inner.  Its quotient formula uses two elements
of the finite-dimensional rational space \(K_\Theta\), so the interpolant is
rational and hence a finite Blaschke product.  This also proves the uniqueness
asserted in the lemma.  The role of such finite-Blaschke extremizers in the
matrix functional-calculus problem is developed by
Crouzeix~\cite{Crouzeix2004} and Li~\cite{Li2020}.
\end{proof}

For a bounded domain \(\Omega\) and a matrix \(A_0\) with
\(\spec(A_0)\subset\Omega\), define
\[
 \psi_\Omega(A_0)=\sup\{\|g(A_0)\|:g\in H^\infty(\Omega),\
                         \|g\|_\Omega\leq1\}.
\]
The value \(g(A_0)\) depends only on finitely many derivatives, those prescribed
by the Jordan blocks.  Montel's theorem therefore makes the supremum a
maximum.  When \(\Omega\) is convex, Lemma~\ref{lem:complexfacts} replaces an
extremizer by a finite Blaschke product composed with a Riemann map, without
changing \(g(A_0)\).

\begin{Lemma}[Strict domain monotonicity]\label{lem:strict-domain}
Let \(\Omega_0\subsetneq\Omega_1\) be bounded convex domains and suppose
\(\spec(A_0)\subset\Omega_0\).  If \(\psi_{\Omega_1}(A_0)>1\), then
\[
 \psi_{\Omega_0}(A_0)>\psi_{\Omega_1}(A_0).
\]
\end{Lemma}

\begin{proof}
Monotonicity in the displayed direction is immediate by restriction.  Suppose
equality held, with common value \(\eta>1\), and choose a norm-one extremizer \(g\)
on \(\Omega_1\).  Its norm on \(\Omega_0\) must be exactly \(1\), since otherwise
rescaling its restriction gives a value larger than \(\eta\).  Thus it is also an
extremizer on \(\Omega_0\).

For either domain, interpolate the finite jet of \(g\) that determines \(g(A_0)\)
with least possible norm.  A norm smaller than \(1\) would contradict
extremality after rescaling.  The uniqueness part of
Lemma~\ref{lem:complexfacts} identifies \(g\) on each domain with a finite
Blaschke product composed with its Riemann map.  In particular, the restriction
of \(g\) to \(\Omega_0\) has a continuous extension to
\(\overline{\Omega_0}\) with modulus one on \(\partial\Omega_0\).

Choose \(z_\star\in\Omega_0\).  Along every ray from \(z_\star\), convexity gives an exit
radius for each domain.  Since the domains differ, on some ray the first exit
occurs strictly earlier; hence there is
\(\xi\in\partial\Omega_0\cap\Omega_1\).  At this interior point of \(\Omega_1\),
continuity of the original holomorphic function and of the boundary extension
gives \(|g(\xi)|=1\).  The maximum-modulus principle makes \(g\) constant,
which would give \(\eta=1\), a contradiction.
\end{proof}

\section{Reduction to an interior-spectrum equality pair}

\begin{Lemma}[Boundary-spectrum induction]\label{lem:induction}
Assume Theorem~\ref{thm:main} has been proved in every size smaller than \(N\).
Let \(A\in M_N(\C)\), let \(K=W(A)\), and suppose \(\psi(A)=2\).  Then either
\(K\) is a nondegenerate closed disk, or
\[
 \spec(A)\subset\Int K.                               \tag{4.1}
\]
\end{Lemma}

\begin{proof}
By Lemma~\ref{lem:foundation}, \(K\) has nonempty interior.  Choose
polynomials \(p_j\) such that
\[
 \|p_j\|_K=1,\qquad \|p_j(A)\|\longrightarrow2.       \tag{4.2}
\]
Suppose \(A\) has a boundary eigenvalue.  Lemma~\ref{lem:boundary-eigen} gives
\(A\simeq A_\partial\oplus A_1\), where \(A_\partial\) is diagonal with
spectrum on \(\partial K\) and \(\spec(A_1)\subset\Int K\).  The
\(A_\partial\)-block of \(p_j(A)\) has norm at most
\(1\), so
\[
 \|p_j(A_1)\|\longrightarrow2.                        \tag{4.3}
\]
In particular \(A_1\) is nonempty.  Since \(W(A_1)\subset K\), (4.3) and
Proposition~\ref{prop:sharp} give \(\psi(A_1)=2\).  Its size is smaller than
\(N\), so induction says that
\[
 K_1:=W(A_1)
\]
is a closed disk of positive radius.

Apply Lemma~\ref{lem:boundary-eigen} once more, now to \(A_1\) relative to
\(K_1\):
\[
 A_1\simeq A_{\partial,1}\oplus A_2,
\]
where \(A_{\partial,1}\) is a finite diagonal block on \(\partial K_1\) and
\(\spec(A_2)\subset\Int K_1\).  The same sequence (4.2) is bounded by \(1\) at
the eigenvalues of \(A_{\partial,1}\), all of which lie in \(K\), so
\(\|p_j(A_2)\|\to2\).  Thus \(A_2\) is nonempty.

Write \(\Lambda_1=\spec(A_{\partial,1})\).  Since
\[
 K_1=W(A_1)=\conv(\Lambda_1\cup W(A_2)),
\]
every extreme point of \(K_1\) belongs to the compact set
\(\Lambda_1\cup W(A_2)\).  Indeed, a finite convex representation of an
extreme point can contain no two distinct points with positive coefficients.
The circle \(\partial K_1\) has
infinitely many extreme points and \(\Lambda_1\) is finite.  Hence
\(\partial K_1\setminus\Lambda_1\subset W(A_2)\); closedness and convexity of
\(W(A_2)\) give
\[
 W(A_2)=K_1.                                           \tag{4.4}
\]

Let \(\Omega_0=\Int K_1\) and \(\Omega_1=\Int K\).  The polynomial sequence and
Proposition~\ref{prop:sharp} show
\[
 \psi_{\Omega_0}(A_2)=\psi_{\Omega_1}(A_2)=2.          \tag{4.5}
\]
Indeed,
\[
 \|p_j\|_{\Omega_0}=\|p_j\|_{K_1}\leq1,
 \qquad
 \|p_j\|_{\Omega_1}=\|p_j\|_K=1,
\]
so the lower bounds follow from (4.2)--(4.4), while the upper bounds follow
from the \(H^\infty\) estimate in Proposition~\ref{prop:sharp}.  If
\(K_1\subsetneq K\), then
\(\Omega_0\subsetneq\Omega_1\), and Lemma~\ref{lem:strict-domain} contradicts
(4.5).  Therefore \(K_1=K\), and \(K\) is the asserted disk.  If this conclusion
does not occur, \(A\) has no boundary eigenvalue, which is (4.1).
\end{proof}

\begin{Lemma}[The holomorphic extremizer]\label{lem:extremizer}
Suppose \(K=W(A)\) has nonempty interior and
\(\spec(A)\subset\Omega=\Int K\).  If \(\psi(A)=2\), there is a nonconstant
\[
 f=\mathcal B\circ\phi\in\mathcal A(K),
\]
where \(\phi:\Omega\to\D\) is conformal and \(\mathcal B\) is a finite
Blaschke product,
such that
\[
 \|f\|_K=1,\qquad |f|=1\text{ on }\partial K,\qquad
 \|f(A)\|=2.                                          \tag{4.6}
\]
\end{Lemma}

\begin{proof}
Normalized polynomials witnessing \(\psi(A)=2\) show
\(\psi_\Omega(A)\geq2\), and Proposition~\ref{prop:sharp} gives the reverse
inequality.  Montel compactness, followed by convergence of the finitely many
spectral jets, gives an \(H^\infty(\Omega)\) maximizer of norm at most one.  Its
norm is exactly one, since otherwise rescaling it would produce a matrix value
of norm greater than two.  Apply the
finite Schur algorithm in Lemma~\ref{lem:complexfacts} to the jet that
determines its matrix value.  It supplies \(\mathcal B\circ\phi\) with the same matrix
value.  The boundary extension of \(\phi\) puts this function in \(\mathcal A(K)\) and
gives boundary modulus \(1\).  It cannot be constant because a constant matrix
has norm at most \(1\).
\end{proof}

\section{From equality to a rational inner function}

\begin{Lemma}[A curved numerical-range arc]\label{lem:arc}
If \(K=W(A)\) has interior, \(\spec(A)\subset\Int K\), and \(K\) is not a
polygon, then \(\partial K\) contains an open real-analytic arc on which every
supporting line exposes one point and the curvature is positive.
Moreover, under \(\spec(A)\subset\Int K\), \(K\) cannot be a polygon.
\end{Lemma}

\begin{proof}
Put
\[
 \mathcal H(\theta)=\Rea(e^{-i\theta}A),\qquad
 h_A(\theta)=\lambda_{\max}(\mathcal H(\theta)).
\]
By Rellich's analytic perturbation theorem~\cite{Rellich1937}, for every
\(\theta_0\) the eigenvalues of the real-analytic Hermitian matrix family
\(\mathcal H(\theta)\) can be numbered, with multiplicity, by real-analytic
functions of \(\theta\) on a neighborhood of \(\theta_0\).  Take
a finite cover of the parameter circle by such intervals.  On each interval,
place two branches in the same class when they agree on a nonempty open
subinterval; the analytic identity theorem then makes them identical on the
whole interval.  Choose one representative of each class and shrink slightly
to a compact subinterval.  The difference of any two distinct representatives
is a nonzero real-analytic function and therefore has only finitely many zeros
there.  Subdividing at the union of these finitely many crossing points fixes
the ordering of all branches; on every resulting open interval their maximum
\(h_A\) is one analytic branch.  A finite cover and a common refinement give
the asserted finite subdivision of the parameter circle.  This is the
finite-dimensional support-function form of the analytic Kippenhahn-curve
description; compare~\cite{PlaumannSinnWeis2021}.

At a differentiability point, the exposed boundary point is
\[
 \sigma(\theta)=e^{i\theta}\bigl(h_A(\theta)+ih_A'(\theta)\bigr),\qquad
 \sigma'(\theta)=ie^{i\theta}\bigl(h_A(\theta)+h_A''(\theta)\bigr). \tag{5.1}
\]
The first identity follows by differentiating the top Rayleigh quotient; the
second follows by differentiation.  Convexity gives
\(h_A+h_A''\geq0\).  If this function vanished identically on every analytic
interval, \(\sigma\) would be constant on each interval and \(K\) would be a
polygon.  Otherwise analyticity supplies a subinterval on which
\(h_A+h_A''>0\).  Formula (5.1) then gives the required regular, positively curved
arc.  Differentiability of the support function makes the exposed face a
singleton there.

It remains to exclude a polygon when the spectrum is interior.  Let
\(\lambda\) be a vertex and let a unit vector \(\upsilon\) realize
\(\upsilon^*A\upsilon=\lambda\).  Two linearly independent supporting normals
\(\nu_1,\nu_2\) at the vertex give
\[
 \Rea(\overline{\nu_j}A)\upsilon
 =\Rea(\overline{\nu_j}\lambda)\upsilon
 \quad(j=1,2).
\]
The two independent real-linear equations imply
\(A\upsilon=\lambda\upsilon\) and
\(A^*\upsilon=\overline\lambda\upsilon\).  Thus \(\lambda\) is a
boundary eigenvalue, contradicting \(\spec(A)\subset\Int K\).
\end{proof}

\begin{Lemma}[Boundary-arc uniqueness]\label{lem:boundary-uniqueness}
Let \(D\subset\C\) be a connected domain and let \(\Gamma_0\subset\partial D\)
be a nonempty open regular real-analytic arc such that, near each point of
\(\Gamma_0\), \(D\) is exactly one of the two local components cut out by the
arc.  If \(g\) is holomorphic
on \(D\), continuous on \(D\cup\Gamma_0\), and its boundary values vanish
arclength-almost everywhere on \(\Gamma_0\), then \(g\equiv0\) on \(D\).
\end{Lemma}

\begin{proof}
Every open subarc has positive arclength, so continuity first makes \(g\) zero
at every point of \(\Gamma_0\).  Near one such point, complexify a regular
real-analytic parametrization of the arc.  Its derivative is nonzero, so the
holomorphic inverse-function theorem gives a biholomorphic coordinate that
sends the arc to a real interval and the local part of \(D\) to one
half-disk.  In that coordinate, extend \(g\) by zero to the opposite
half-disk.  The extension is continuous; subdividing any triangle along the
diameter and applying Cauchy's theorem on the pieces proves Morera's
criterion.  The extension is holomorphic and zero on an open half-disk, hence
zero throughout the disk by the identity theorem.  Thus \(g\) vanishes on a
nonempty open subset of \(D\), and connectedness gives the conclusion.
\end{proof}

\begin{Proposition}[Rational collapse]\label{prop:rational}
Under the hypotheses and notation of Lemma~\ref{lem:extremizer}, let \(x,y\)
be the equality vectors from Proposition~\ref{prop:equality-data}.  Then the
function \(f\) is rational on \(\Omega\).  More precisely, for
\(z\notin\spec(A)\), set
\[
 a(z)=x^*(zI-A)^{-1}x,\qquad c(z)=y^*(zI-A)^{-1}x,
\]
and regard \(a,c\) thereafter as their meromorphic adjugate continuations.
Then
\[
 c(z)f(z)=2a(z)                                       \tag{5.2}
\]
as an identity of meromorphic functions on \(\Omega\).  Equivalently,
\(f=2a/c\) wherever the quotient is defined, and then everywhere after
removable cancellation.
After cancellation, \(f=U/V\) for coprime polynomials with
\[
 \deg U>\deg V.                                       \tag{5.3}
\]
If \(\adj\) denotes the classical adjugate and
\[
 D_A(z)=x^*\adj(zI-A)x,\qquad C_A(z)=y^*\adj(zI-A)x,
\]
then the reduction can be chosen so that
\[
 U(z)C_A(z)=2D_A(z)V(z).
\]
\end{Proposition}

\begin{proof}
Choose the arc from Lemma~\ref{lem:arc}; it avoids the finite spectrum.  For
almost every point \(\sigma\) of that arc, let \(\mathfrak n\) be its outward
normal, \(R=(\sigma I-A)^{-1}\), and
\[
 \Delta_\sigma=\Rea\bigl(\overline{\mathfrak n}(\sigma I-A)\bigr)\succeq0.
\]
Equations (3.5) and (3.9) imply
\[
 \Delta_\sigma R\bigl(y-f(\sigma)x\bigr)=0.           \tag{5.4}
\]
Indeed, squaring the norm in (3.9) and using
\(\mathcal P(\sigma)=\pi^{-1}R^*\Delta_\sigma R\) first gives
\(\Delta_\sigma^{1/2}R(y-f(\sigma)x)=0\), and hence (5.4).
The vector
\(r_\sigma=R(y-f(\sigma)x)\) is nonzero: \(R\) is invertible, while the
orthonormality of \(x,y\) makes \(y-f(\sigma)x\neq0\).  It is therefore an
eigenvector of \(\Rea(\overline{\mathfrak n}A)\) for its maximal eigenvalue.
The exposed point is unique, so
\(r_\sigma^*Ar_\sigma=\sigma\|r_\sigma\|^2\).  Since
\((\sigma I-A)r_\sigma=y-f(\sigma)x\), taking the inner product with
\(r_\sigma\) and conjugating gives
\[
 \bigl(y-f(\sigma)x\bigr)^*
 R\bigl(y-f(\sigma)x\bigr)=0.                         \tag{5.5}
\]

For \(z\notin\spec(A)\), put \(R_z=(zI-A)^{-1}\).  This resolvent commutes
with \(T=f(A)\), and relations (3.8) give
\[
 x^*R_z y=\tfrac12x^*R_zTx=\tfrac12x^*TR_zx=0,
 \qquad
 y^*R_z y=\tfrac12y^*R_zTx=\tfrac12y^*TR_zx=a(z).   \tag{5.6}
\]
Expanding (5.5), using \(|f(\sigma)|=1\) and (5.6), yields
\[
 2a(\sigma)-f(\sigma)c(\sigma)=0                     \tag{5.7}
\]
almost everywhere on the arc.  The punctured domain
\(D=\Omega\setminus\spec(A)\) is connected, and \(cf-2a\) is holomorphic on
\(D\) and continuous up to the chosen arc.  Convexity supplies the local
one-sidedness required by Lemma~\ref{lem:boundary-uniqueness}; that lemma gives
\(cf-2a=0\) on \(D\), hence as a meromorphic identity on \(\Omega\).  Since
\(a(z)=z^{-1}+O(z^{-2})\) at infinity, \(a\) is not the
zero rational function; the identity therefore also shows that \(c\) is not
identically zero.  This proves (5.2).

The adjugate formulas are
\[
 a=\frac{D_A}{\det(zI-A)},\qquad
 c=\frac{C_A}{\det(zI-A)}.
\]
The polynomial \(D_A\) has degree exactly \(N-1\), with leading coefficient
\(1\).  Since \(x^*y=0\), the degree of \(C_A\) is at most \(N-2\).  Let
\(G\) be a greatest common divisor of \(C_A\) and \(D_A\), and choose
\[
 U=\frac{2D_A}{G},\qquad V=\frac{C_A}{G}.
\]
Then \(U,V\) are coprime, \(f=U/V\), and
\(UC_A=2D_AV\).  Canceling \(G\) preserves the strict degree difference,
proving (5.3).
\end{proof}

\section{The full lemniscate is the boundary}

\begin{Lemma}[Zeros of the transfer function]\label{lem:zeros}
Let \(f=U/V\) be the reduced rational function from
Proposition~\ref{prop:rational}.  Then every finite zero of \(U\) lies in
\(\Omega\), every zero of \(V\) lies outside \(K\), and \(f(\infty)=\infty\).
\end{Lemma}

\begin{proof}
For \(z\notin K\), put \(R=(zI-A)^{-1}\).  Then \(a(z)\neq0\).
Indeed, if \(a(z)=x^*Rx=0\) and \(v=Rx\), then \(v\neq0\) and
\((zI-A)v=x\).  Since \(x^*v=a(z)=0\), conjugation gives \(v^*x=0\), and hence
\[
 v^*(zI-A)v=v^*x=0,
 \qquad\text{hence}\qquad
 \frac{v^*Av}{\|v\|^2}=z,
\]
which would put \(z\) in \(W(A)=K\), a contradiction.

By the construction in Proposition~\ref{prop:rational}, \(U=2D_A/G\), so every
zero of \(U\) is a zero of \(D_A\).  Since \(\det(zI-A)\neq0\) off \(K\) and
\(a=D_A/\det(zI-A)\) does not vanish there, \(U\) has no zero off \(K\).  Nor
can \(V\) vanish at a boundary point: by
coprimality, \(U/V\) would then be unbounded along approaches from \(\Omega\),
where it equals the original function \(f\in\mathcal A(K)\).  Thus \(U/V\)
extends continuously to \(\partial K\), agrees there with \(f\), and has
modulus one.  In particular, \(U\) has no boundary zero, so all its zeros lie
in \(\Omega\).  A reduced pole cannot lie in \(\Omega\), where \(f\) is
holomorphic.  Hence all poles lie outside \(K\).
Finally (5.3) says exactly that \(f(\infty)=\infty\).
\end{proof}

\begin{Proposition}[Full-level identity]\label{prop:full-level}
For the reduced rational extremizer \(f=U/V\),
\[
 \{z\in\C:V(z)\neq0,\ |U(z)/V(z)|<1\}=\Omega,
 \qquad
 \{z\in\C:V(z)\neq0,\ |U(z)/V(z)|=1\}=\partial K.   \tag{6.1}
\]
Moreover \(f'\) does not vanish on \(\partial K\), and \(\partial K\) is a
smooth real-analytic strictly convex Jordan curve.
\end{Proposition}

\begin{proof}
Put
\[
 \mathcal O:=\{z\in\C:V(z)\neq0,\ |U(z)/V(z)|<1\}.
\]
No pole belongs to \(\overline{\mathcal O}\), because \(|f|>1\) in a
punctured neighborhood of each pole.  If \(z\in\partial\mathcal O\), continuity
and approximation from \(\mathcal O\) give \(|f(z)|\leq1\); strict inequality
would put a neighborhood of \(z\) in \(\mathcal O\).  Therefore
\[
 \partial\mathcal O\subset
 \{z\in\C:V(z)\neq0,\ |U(z)/V(z)|=1\}.              \tag{6.2}
\]
Now \(\Omega\subset\mathcal O\), while \(\partial K\) is disjoint from
\(\mathcal O\).  Hence the connected set \(\Omega\) is both open and closed
relative to \(\mathcal O\), and is one of its components.

Every component \(\mathcal O_0\) is bounded because \(f(\infty)=\infty\), and
it contains a zero of \(f\).  Otherwise its compact closure would avoid all
zeros of \(f\): any zero in the closure lies in \(\mathcal O\) and hence in
\(\mathcal O_0\).  Thus \(1/f\) would be holomorphic near
\(\overline{\mathcal O_0}\), have modulus one on the boundary by (6.2), and
have modulus greater than one inside, contrary to the maximum-modulus
principle.  Lemma~\ref{lem:zeros} puts every zero in \(\Omega\), so
\(\mathcal O\) has no other component.  This proves the first identity in
(6.1).

Conversely, if \(V(z_0)\neq0\) and
\(|f(z_0)|=1\), the open mapping theorem says that every neighborhood of
\(z_0\) contains points where \(|f|<1\).  Thus
\(z_0\in\partial\mathcal O\), and (6.2) proves the second identity in (6.1).

Suppose that \(f'(z_0)=0\) at a level point.  Then, in a neighborhood of
\(z_0\),
\[
 f(z)-f(z_0)=(z-z_0)^kq(z),\qquad k\geq2,\quad q(z_0)\neq0.
\]
Choose \(\alpha\in\C\) so that
\(\Rea\bigl(\overline{f(z_0)}\alpha^kq(z_0)\bigr)<0\), and let
\(\zeta=e^{2\pi i/k}\).  For all sufficiently small \(t>0\), the points
\(z_j=z_0+t\alpha\zeta^j\), \(0\leq j<k\), are not poles and satisfy
\[
 |f(z_j)|^2
 =1+2t^k\Rea\bigl(\overline{f(z_0)}\alpha^kq(z_0)\bigr)
   +O(t^{k+1})<1.
\]
The first identity in (6.1) puts every \(z_j\) in \(\Omega\).  But
\[
 \frac1k\sum_{j=0}^{k-1}z_j=z_0
\]
because \(\sum_{j=0}^{k-1}\zeta^j=0\).  Convexity of \(\Omega\) would then put
the boundary point \(z_0\) in \(\Omega\), a contradiction.  Hence
\(f'\neq0\) on the level, and the implicit-function theorem makes the boundary
real analytic and smooth.

Finally suppose that the boundary contained a nontrivial segment of a real
line \(z=z_0+t\xi\), \(t\in\R\), \(\xi\neq0\).  The expression
\[
 |U(z_0+t\xi)|^2-|V(z_0+t\xi)|^2
\]
is a polynomial in the real variable \(t\).  Its vanishing on an interval
would make it vanish for every real \(t\).  If \(V(z_0+t\xi)=0\) for some real
\(t\), the same identity would force \(U(z_0+t\xi)=0\), contrary to
coprimality.  The full-level identity would therefore put the entire
unbounded line in \(\partial K\), contradicting compactness.  Thus the boundary
has no nontrivial segment, which for a compact convex body is exactly strict
convexity.
\end{proof}

\section{The algebraic tangent argument}

The Hermitian determinant below is the line-coordinate polynomial introduced
by Kippenhahn~\cite{Kippenhahn2008}; modern projective-dual
accounts appear in~\cite{PlaumannSinnWeis2021}.  Throughout this section,
primal projective coordinates are ordered as \([Z:X:Y]\), dual coordinates as
\([s:u:v]\), and incidence is
\[
 sZ+uX+vY=0.
\]
Thus the affine point \(z=X+iY\) is \([1:X:Y]\).

\begin{Proposition}[Generic hyperbolic contact count]
\label{prop:generic-contact}
Let \(\mathscr C\subset\mathbb P^2_\C\) be an irreducible projective curve
defined over \(\R\), and suppose its entire real locus is a smooth strictly
convex oval in the affine chart \(Z=1\).  Let
\(\mathscr Q(s,u,v)\) be a real irreducible homogeneous equation of the dual
curve \(\mathscr C^*\).  If
\[
 \mathscr Q(1,0,0)\neq0
 \quad\text{and}\quad
 s\longmapsto\mathscr Q(s,u,v)
 \text{ has only real zeros for all }(u,v)\in\R^2,
\]
then \(\deg\mathscr Q=2\).
\end{Proposition}

\begin{proof}
Write \(m=\deg\mathscr Q\).  The normalization, conormal, and biduality theory
of projective plane curves in characteristic zero gives the following generic
contact description; see Tevelev~\cite{Tevelev2003}.
\begin{enumerate}
\item Outside a finite subset \(B\subset\mathscr C^*\), a smooth dual point
has one smooth contact point on \(\mathscr C\), and the tangent line to
\(\mathscr C^*\) at that dual point recovers the contact point.  This is the
generic one-to-one Gauss correspondence supplied by biduality; \(B\) absorbs
the singular, branch, and ramification loci.
\item Since \([1:0:0]\notin\mathscr C^*\), projection from that point,
\[
 \pi:\mathscr C^*\longrightarrow\mathbb P^1,
 \qquad [s:u:v]\longmapsto[u:v],
\]
is everywhere defined and has generic degree \(m\): a generic fiber is the
intersection with a line through \([1:0:0]\), and B\'ezout gives \(m\) points.
\item In characteristic zero this projection is separable.  Equivalently, the
discriminant of \(\mathscr Q(s,u,v)\), viewed as a binary homogeneous
polynomial in \((u,v)\), is not identically zero.  The branch values of \(\pi\)
and the projection of \(B\) therefore form a finite exceptional subset of
\(\mathbb P^1\).
\item Choose a real \([u:v]\) outside that set.  Because the coefficient of
\(s^m\) is \(\mathscr Q(1,0,0)\neq0\), specialization does not lower the
degree.  Hence there are \(m\) distinct roots \(s_j\), all real by hypothesis,
and \(\ell_j=[s_j:u:v]\) are smooth real points of \(\mathscr C^*\setminus B\).
\item Biduality and the chosen coordinate order identify the unique contact
point with
\[
 [Z:X:Y]=
 [\mathscr Q_s(\ell_j):\mathscr Q_u(\ell_j):
   \mathscr Q_v(\ell_j)].
\]
The polynomial and \(\ell_j\) are real, so this point is real and therefore
lies on the given oval.
\item A fixed unoriented normal \([u:v]\) has exactly two tangent lines to a
smooth strictly convex oval: the unique support lines at the maximum and
minimum of \(uX+vY\).  Thus the \(m\) lines just found give \(m\leq2\).
Conversely, both support lines are points of \(\mathscr C^*\), so they are two
distinct roots and \(m\geq2\).  Hence \(m=2\).
\end{enumerate}
\end{proof}

\begin{Proposition}[Circle criterion]\label{prop:circle}
Let \(K\) be a compact convex body whose boundary is the full level set
\[
 \partial K=\{z\in\C:V(z)\neq0,\ |U(z)/V(z)|=1\},    \tag{7.1}
\]
where \(U,V\) are coprime polynomials, \(\deg U>\deg V\), \(K\) is strictly
convex, and \((U/V)'\) does not vanish on \(\partial K\).  Writing
\(z=X+iY\), suppose there are Hermitian matrices \(\mathsf H,\mathsf J\) such
that, in the real dual coordinates
\([s:u:v]\) associated with the line \(s+uX+vY=0\), every tangent line to an
open boundary arc belongs to the projective curve
\[
 \det(sI+u\mathsf H+v\mathsf J)=0.                   \tag{7.2}
\]
Then \(\partial K\) is a circle.
\end{Proposition}

\begin{proof}
Put \(f=U/V\).  Write \(z=X+iY\), \(w=X-iY\), and
\(S^\#(w)=\overline{S(\overline w)}\) for a polynomial \(S\).
For a homogenization \(S_h(w,Z)\), the notation \(S_h^\#\) means
coefficientwise conjugation, with \(w,Z\) left unchanged.
Put \(\mathbf z=(Z,X,Y)^{\mathsf T}\) and
\(\boldsymbol\ell=(s,u,v)^{\mathsf T}\), in the coordinate order fixed above.
Let \(d=\deg U\), \(e=\deg V<d\), and homogenize the real lemniscate polynomial
to degree \(2d\):
\[
 \mathscr F(Z,X,Y)=
 U_h(X+iY,Z)U_h^\#(X-iY,Z)
 -Z^{2(d-e)}V_h(X+iY,Z)V_h^\#(X-iY,Z).               \tag{7.3}
\]
Here \(U_h,V_h\) are the usual homogenizations to their own degrees.  At
\(Z=1\), use the Wirtinger derivative
\(\partial_z=\tfrac12(\partial_X-i\partial_Y)\).  Then
\[
 \mathscr F=|V|^2(|f|^2-1),\qquad
 \partial_z\mathscr F=|V|^2 f'\overline f
 \quad\text{on }|f|=1.
\]
Here \(V\neq0\) on the level and \(f'\neq0\) there by hypothesis, so the
boundary lies in the regular locus of \(\mathscr F=0\).

\medskip\noindent\emph{Selection of the real component.}
Factor \(\mathscr F\) over \(\C\).  At a regular zero exactly one irreducible
factor occurrence vanishes: two vanishing factors, or a repeated vanishing
factor, would make every first derivative of \(\mathscr F\) zero.  The choice
of that factor is locally constant along the boundary, because all other
factors remain nonzero nearby.  The boundary is connected, so one factor
\(\mathscr P\) vanishes on all of it; let \(\mathscr C=V(\mathscr P)\).
Since \(\mathscr F\) has real coefficients, the coefficientwise conjugate
\(\overline{\mathscr P}\) is also a factor and vanishes at each real boundary
point.  Local uniqueness makes \(\overline{\mathscr P}\) an associate of
\(\mathscr P\); rescaling \(\mathscr P\) therefore gives real coefficients.
Thus \(\mathscr C\) is an irreducible real projective curve, and it is smooth
along the boundary.

\medskip\noindent\emph{The entire real locus.}
Every real affine point of \(\mathscr C\) is a real zero of (7.3).  If
\(V(z)=0\) there, (7.3) also forces \(U(z)=0\), contradicting coprimality.
Hence \(V(z)\neq0\), and (7.1) puts the point on \(\partial K\).  The reverse
inclusion holds by construction of \(\mathscr C\), so
\[
 \mathscr C(\R)\cap\{Z\neq0\}
 =\{[1:X:Y]:X+iY\in\partial K\}.
\]
There are no additional real points at infinity.  Indeed, if
\(\operatorname{lc}(U)\) denotes the nonzero leading coefficient of \(U\),
\[
 \mathscr F(0,X,Y)=|\operatorname{lc}(U)|^2
                    (X^2+Y^2)^d,                    \tag{7.4}
\]
and \(X^2+Y^2>0\) for every nonzero real projective pair \([X:Y]\).  Therefore
\[
 \mathscr C(\R)=\{[1:X:Y]:X+iY\in\partial K\},
\]
exactly one smooth strictly convex affine oval.

Let \(\mathscr C^*\) be the projective dual curve, the Zariski closure of the
tangent lines at smooth points of \(\mathscr C\), and let
\(\mathscr Q(s,u,v)\) be its irreducible homogeneous equation, chosen with
real coefficients.  Tangency on an open real arc and (7.2) imply
\[
 \mathscr Q(s,u,v)\ \text{divides}
 \det(sI+u\mathsf H+v\mathsf J).                     \tag{7.5}
\]
Indeed, the Gauss image of that arc is nonconstant: a constant tangent line
would contain a nontrivial boundary arc, contrary to strict convexity.  Its
image is therefore infinite and hence Zariski dense in the irreducible curve
\(\mathscr C^*\).  The determinant vanishes on that dense set and thus on all
of \(\mathscr C^*\).  The homogeneous prime ideal of a plane curve is the
principal ideal \((\mathscr Q)\), proving (7.5).  Since \(\mathscr C\) is real,
its dual is invariant under conjugation; as in the primal factor argument,
\(\mathscr Q\) can be rescaled to have real coefficients.

Set \(\mathbf e=[1:0:0]\).  The determinant in (7.5) equals \(1\) at
\(\mathbf e\), so \(\mathscr Q(1,0,0)\neq0\).  If
\(m=\deg\mathscr Q\), homogeneity gives
\[
 [s^m]\mathscr Q(s,u,v)=\mathscr Q(1,0,0)\neq0.
\]
Thus specialization never lowers the degree in \(s\).  Specializing the
polynomial divisibility (7.5) shows that every zero of
\(\mathscr Q(s,u,v)\) is a zero of the Hermitian determinant.  For real
\(u,v\), the latter zeros are the negatives of the eigenvalues of
\(u\mathsf H+v\mathsf J\), and are therefore real.  Hence \(\mathscr Q\) is
hyperbolic with respect to \(\mathbf e\) in the elementary real-rootedness
sense~\cite{Garding1959,HeltonVinnikov2007}.

The exact real-locus identity above and
Proposition~\ref{prop:generic-contact} now give
\[
 \deg\mathscr Q=2.                                    \tag{7.6}
\]
Thus \(\mathscr C^*\) is an irreducible conic and is nonsingular.  Writing its
equation as
\(\boldsymbol\ell^{\mathsf T}\Sigma\boldsymbol\ell=0\), with an invertible
symmetric \(3\times3\) matrix \(\Sigma\), its tangent at
\(\boldsymbol\ell\) corresponds to the primal point
\(\Sigma\boldsymbol\ell\).  The envelope of those tangents is
\(\mathbf z^{\mathsf T}\Sigma^{-1}\mathbf z=0\), also a conic.  Since it
contains the generic bidual contact points, a Zariski-dense subset of
\(\mathscr C\), it is \(\mathscr C\).

Finally, because the defining polynomial of \(\mathscr C\) divides
\(\mathscr F\), every point of this conic at infinity must be one of
\[
 [0:1:i],\qquad[0:1:-i].
\]
The line \(Z=0\) is not a component of \(\mathscr C\): the curve is
irreducible and contains its affine oval.  B\'ezout's theorem therefore gives
total intersection multiplicity two with the line at infinity.  Complex
conjugation exchanges the two
displayed, distinct circular points and preserves multiplicity, so each
occurs once.  The real homogeneous equation of the conic consequently has the
form
\[
 q_0(X^2+Y^2)+q_1ZX+q_2ZY+q_3Z^2=0,
 \qquad q_0\neq0.
\]
Completing squares shows that its nonempty compact real oval is a circle.
\end{proof}

\section{Proof of the main theorem}

\begin{proof}[Proof of Theorem~\ref{thm:main}]
We use induction on \(N\).  For \(N=1\), and more generally whenever the
numerical range has empty interior,
Lemma~\ref{lem:foundation} gives \(\psi(A)=1\), so the hypothesis never occurs.

Assume the theorem holds for all sizes smaller than \(N\), and let
\(A\in M_N(\C)\) satisfy
\(\psi(A)=2\).  Put \(K=W(A)\).  It has nonempty interior.  By
Lemma~\ref{lem:induction}, either \(K\) is already a nondegenerate disk or
\(\spec(A)\subset\Omega=\Int K\).  Consider the second case.

Choose \(f\in\mathcal A(K)\) as in Lemma~\ref{lem:extremizer}, and let \(x,y\)
be the equality vectors from Proposition~\ref{prop:equality-data}.  Since an
interior-spectrum numerical range cannot be a polygon, Lemma~\ref{lem:arc}
provides a curved analytic boundary arc.  On that arc,
Proposition~\ref{prop:rational} yields
\[
 f(z)=\frac{U(z)}{V(z)},\qquad \deg U>\deg V,
\]
after cancellation of the meromorphic transfer identity (5.2).

Lemma~\ref{lem:zeros} puts every zero of \(U\) in \(\Omega\), every pole
outside \(K\), and gives \(f(\infty)=\infty\).
Proposition~\ref{prop:full-level} then gives
\[
 \{z\in\C:V(z)\neq0,\ |U(z)/V(z)|<1\}=\Omega,
 \qquad
 \{z\in\C:V(z)\neq0,\ |U(z)/V(z)|=1\}=\partial K,
\]
and shows that this boundary is a smooth strictly convex rational lemniscate.

Put \(H_A=\Rea A\) and \(J_A=\Ima A\).  Writing \(z=X+iY\), every supporting
line \(s+uX+vY=0\), with \(s,u,v\in\R\), can be oriented so that
\(sI+uH_A+vJ_A\succeq0\).  A unit vector attaining the point of support has
zero quadratic form for this positive semidefinite matrix and therefore lies
in its kernel.  Hence
\[
 \det(sI+uH_A+vJ_A)=0.
\]
All hypotheses of Proposition~\ref{prop:circle} now hold.  It follows that
\(\partial K\) is a circle.  Compactness and convexity make \(K\) the filled
closed disk bounded by that circle.

The disk has positive radius because \(K\) has nonempty planar interior.
Writing its center and radius as \(\gamma\) and \(\rho\) gives
\[
 W(A)=\{z:|z-\gamma|\leq\rho\},\qquad \rho>0,
\]
which completes the induction and the proof.
\end{proof}

\funding{This research received no external funding.}
\acknowledgments{The author used Codex with GPT-5.6 Sol Ultra for proof development,
Lean formalization, and language editing, reviewed all output, and takes full
responsibility for the publication.}
\dataavailability{The source archive contains the theorem-matched Lean
development.  The audited \texttt{.lean} sources and toolchain files
correspond to Git commit
\texttt{\seqsplit{5af901c72a5148f782cc91551f76127ec93a98c5}} and use
Lean~4.33.1 with mathlib~v4.33.1.  The file
\texttt{Lean/FORMALIZATION.md} records the formalization coverage,
dependencies, and reproducible fresh-checkout commands.}
\conflictsofinterest{The author declares no conflicts of interest.}

\reftitle{References}

\begin{thebibliography}{99}

\bibitem[Crouzeix(2004)]{Crouzeix2004}
Crouzeix, M. Bounds for analytical functions of matrices.
\emph{Integral Equations Oper. Theory} \textbf{2004}, \emph{48}, 461--477.
\url{https://doi.org/10.1007/s00020-002-1188-6}.

\bibitem[Crouzeix(2007)]{Crouzeix2007}
Crouzeix, M. Numerical range and functional calculus in Hilbert space.
\emph{J. Funct. Anal.} \textbf{2007}, \emph{244}, 668--690.
\url{https://doi.org/10.1016/j.jfa.2006.10.013}.

\bibitem[Okubo and Ando(1975)]{OkuboAndo1975}
Okubo, K.; Ando, T. Constants related to operators of class \(C_{\rho}\).
\emph{Manuscripta Math.} \textbf{1975}, \emph{16}, 385--394.
\url{https://doi.org/10.1007/BF01323467}.

\bibitem[Delyon and Delyon(1999)]{DelyonDelyon1999}
Delyon, B.; Delyon, F. Generalization of von Neumann's spectral sets and
integral representation of operators. \emph{Bull. Soc. Math. France}
\textbf{1999}, \emph{127}, 25--41.
\url{https://doi.org/10.24033/bsmf.2340}.

\bibitem[Crouzeix and Palencia(2017)]{CrouzeixPalencia2017}
Crouzeix, M.; Palencia, C. The numerical range is a
\((1+\sqrt{2})\)-spectral set. \emph{SIAM J. Matrix Anal. Appl.}
\textbf{2017}, \emph{38}, 649--655.
\url{https://doi.org/10.1137/17M1116672}.

\bibitem[Jin(2026)]{Jin2026}
Jin, S. The numerical range is a \(2\)-spectral set.
\emph{Preprints} \textbf{2026}, 202607.1919.v4.
\url{https://doi.org/10.20944/preprints202607.1919.v4}.

\bibitem[Lorist and Schwenninger(2026)]{LoristSchwenninger2026}
Lorist, E.; Schwenninger, F.L. A solution to Crouzeix's conjecture.
\emph{arXiv} \textbf{2026}, arXiv:2608.03841v2, 7~pp.
\url{https://doi.org/10.48550/arXiv.2608.03841}.

\bibitem[{\AA}hag et al.(2026)]{AhagEtAl2026}
{\AA}hag, P.; Czy\.z, R.; Per\"al\"a, A.; Virtanen, J.
Square functions and the complete Crouzeix conjecture in dimension three.
\emph{arXiv} \textbf{2026}, arXiv:2608.27346v1.
\url{https://doi.org/10.48550/arXiv.2608.27346}.

\bibitem[Li(2020)]{Li2020}
Li, K. On the uniqueness of functions that maximize the Crouzeix ratio.
\emph{Linear Algebra Appl.} \textbf{2020}, \emph{599}, 105--120.
\url{https://doi.org/10.1016/j.laa.2020.03.045}.

\bibitem[Greenbaum and Overton(2018)]{GreenbaumOverton2018}
Greenbaum, A.; Overton, M.L. Numerical investigation of Crouzeix's
conjecture. \emph{Linear Algebra Appl.} \textbf{2018}, \emph{542}, 225--245.
\url{https://doi.org/10.1016/j.laa.2017.04.035}.

\bibitem[Schwenninger and de Vries(2025)]{SchwenningerDeVries2025}
Schwenninger, F.L.; de Vries, J. The double-layer potential for spectral
constants revisited. \emph{Integral Equations Oper. Theory} \textbf{2025},
\emph{97}, Article~13, 22~pp.
\url{https://doi.org/10.1007/s00020-025-02800-2}.

\bibitem[Malman et al.(2025)]{MalmanEtAl2025}
Malman, B.; Mashreghi, J.; O'Loughlin, R.; Ransford, T. Double-layer
potentials, configuration constants, and applications to numerical ranges.
\emph{Int. Math. Res. Not. IMRN} \textbf{2025}, 2025, Paper No.~rnaf084,
34~pp. \url{https://doi.org/10.1093/imrn/rnaf084}.

\bibitem[Sarason(1967)]{Sarason1967}
Sarason, D. Generalized interpolation in \(H^{\infty}\).
\emph{Trans. Am. Math. Soc.} \textbf{1967}, \emph{127}, 179--203.
\url{https://doi.org/10.1090/S0002-9947-1967-0208383-8}.

\bibitem[Bultheel and Lasarow(2008)]{BultheelLasarow2008}
Bultheel, A.; Lasarow, A. Solution of a multiple Nevanlinna--Pick problem
for Schur functions via orthogonal rational functions.
\emph{Analysis} \textbf{2008}, \emph{28}, 95--123.
\url{https://doi.org/10.1524/anly.2008.0904}.

\bibitem[Rellich(1937)]{Rellich1937}
Rellich, F. St\"orungstheorie der Spektralzerlegung. I. Mitteilung.
Analytische St\"orung der isolierten Punkteigenwerte eines beschr\"ankten
Operators. \emph{Math. Ann.} \textbf{1937}, \emph{113}, 600--619.
\url{https://doi.org/10.1007/BF01571652}.

\bibitem[Kippenhahn(2008)]{Kippenhahn2008}
Kippenhahn, R. On the numerical range of a matrix; Zachlin, P.F.;
Hochstenbach, M.E., Translators. \emph{Linear Multilinear Algebra}
\textbf{2008}, \emph{56}, 185--225.
\url{https://doi.org/10.1080/03081080701553768}.

\bibitem[Plaumann et al.(2021)]{PlaumannSinnWeis2021}
Plaumann, D.; Sinn, R.; Weis, S. Kippenhahn's theorem for joint numerical
ranges and quantum states. \emph{SIAM J. Appl. Algebra Geom.}
\textbf{2021}, \emph{5}, 86--113.
\url{https://doi.org/10.1137/19M1286578}.

\bibitem[G{\aa}rding(1959)]{Garding1959}
G{\aa}rding, L. An inequality for hyperbolic polynomials.
\emph{J. Math. Mech.} \textbf{1959}, \emph{8}, 957--965.
\url{https://doi.org/10.1512/iumj.1959.8.58061}.

\bibitem[Helton and Vinnikov(2007)]{HeltonVinnikov2007}
Helton, J.W.; Vinnikov, V. Linear matrix inequality representation of sets.
\emph{Commun. Pure Appl. Math.} \textbf{2007}, \emph{60}, 654--674.
\url{https://doi.org/10.1002/cpa.20155}.

\bibitem[Tevelev(2003)]{Tevelev2003}
Tevelev, E.A. Projectively dual varieties.
\emph{J. Math. Sci.} \textbf{2003}, \emph{117}, 4585--4732.
\url{https://doi.org/10.1023/A:1025366207448}.

\end{thebibliography}

\end{document}
